\chapter{利率的传导链：从政策利率到你的房贷}
\chaptermeta{09}

\begin{ledetext}
央行在最上游拧了一下阀门。这股水要过七道闸才能流到陈宁每月 15 号被划走的那笔钱上，每一道都可能放大它、拖慢它，或者把它堵死。
\end{ledetext}

\section{从 1.4\% 到 13471 元，中间隔着七道闸}

陈宁的房贷 300 万，30 年，等额本息，执行利率 3.5\%，月供 13471 元。他和妻子加起来到手 3 万 4 出头，月供占了四成。

央行公布的那个数字是 1.4\%。陈宁账单上的数字是 3.5\%。

这两个数字之间隔着一整条管道。

\begin{flowlist}
  \flowitem{01}{政策利率}{中国是 7 天期逆回购利率，截至 2026 年 8 月为 1.4\%；美国是联邦基金利率目标区间 3.50\% 到 3.75\%。全链条的锚。\emph{信号在这里最强、最干净。}}
  \flowitem{02}{银行间货币市场}{中国看 \term{DR007}{存款类机构质押式回购利率}\index{095@DR007}，也就是银行之间拿国债作抵押借 7 天钱的实际成交价；美国看 \term{SOFR}{担保隔夜融资利率}\index{096@SOFR}。\emph{季末、缴税、大额国债发行都会让它偏离政策利率。}}
  \flowitem{03}{同业存单与短期国债}{银行发同业存单筹 3 个月、6 个月的钱，价格直接挂在货币市场利率上。\emph{信用和期限的因素从这里开始掺进来。}}
  \flowitem{04}{收益率曲线}{从 1 个月到 30 年，市场报出一整条价格。\emph{整条链上信号最容易变形的一环，因为长端不归央行管。}}
  \flowitem{05}{银行的负债成本}{银行的钱从存款、同业存单、央行借款三处来。\hlmark{存款利率是刚性的。}没人会因为央行降息就接受活期变负数，银行也不敢把定期利率砍到没人肯存。\emph{信号在这里被结结实实卡掉一大截。}}
  \flowitem{06}{贷款定价}{中国走 LPR 加点，美国各家银行走自己的模型。\emph{银行的资本状况和风险偏好，决定它愿意让出多少。}}
  \flowitem{07}{陈宁的月供}{还得等到他的重定价日。\emph{信号走到这里，可能比出发时晚了整整一年。}}
\end{flowlist}

第 7 章那个电网的比喻，在这里能接着往下推。高压电送到你家插座要走上百公里，每一段线路都有损耗，每一台变压器都要吃掉一点电压。这条链最大的那处损耗在第 5 环：银行的存款成本降不下来，贷款利率就让不出去。

"央行降息 25 个基点"和"陈宁月供少 415 块"之间不是等号，是一段漫长而且可能中断的旅程。

\section{收益率曲线：一张能看出天气的图}

第 4 环那条\term{收益率曲线}{yield curve}\index{053@收益率曲线}值得单独讲，它是整个金融市场信息密度最高的一张图。

画法极简单。横轴是期限（1 个月、1 年、5 年、10 年、30 年），纵轴是对应的国债收益率，把点连起来就是一条线。

\textbf{正常形态向上倾斜。}把钱借出去 10 年，要承担的不确定性远多于借出去一天：通胀可能失控，你可能急着用钱，发行人可能变差。市场为这份不确定性索取的额外补偿，叫\term{期限溢价}{term premium}\index{042@期限溢价}。

曲线变平，说明市场认为利率不会再往上走了。

短端反过来高过长端，叫倒挂。这个形状很刺眼：\hlmark{市场在说，我宁可锁定 10 年拿一个更低的利率，也不愿意赌未来还有这么高的短期利率。}翻译过来就是，大家预期央行将来会大幅降息，而央行大幅降息通常只有一个原因。

倒挂在历史上是可靠的衰退先行指标，还有第二个更机械的理由：\textbf{银行的商业模式是借短贷长。}短端高过长端，银行的息差被直接压扁，放贷动力下降，信用收缩，然后经济真的坏了。它不只是预言，它本身就参与制造了那个结果。

\begin{deepbox}{进阶：可以跳过 · 这个指标失灵了吗}
2022 到 2023 年，美国出现了史上持续最久的倒挂之一，几乎所有人都在等衰退。衰退迟迟没来，经济反而维持了相当的韧性。

解释有三种，都讲得通，都没被证实。一是 QE 长期压低了期限溢价，曲线形状的含义已经和历史上不一样了。二是疫情后居民和企业锁定了大量长期低息负债，资产负债表异常健康，对短期利率不敏感。三是财政持续扩张，对冲掉了货币紧缩。

\textbf{这个指标是不是已经失效，目前没人知道答案。}本书的建议是把倒挂当成一个值得警惕的信号，不是一个可以照着下注的时钟。任何"指标一响就该卖出"的说法都不构成投资建议，历史样本只有寥寥几次，根本撑不起那种确定性。

\end{deepbox}

\section{LPR：一条你能亲眼看见的链子}

中国的贷款定价机制透明得罕见，每一环都能看清。

\begin{chainflow}
\chainnode{7 天逆回购 1.4\%}\chainarrow \chainnode{18 家报价行加点报价}\chainarrow \chainnode{同业拆借中心每月 20 日发布 LPR}\chainarrow \chainnode{房贷 = LPR ± 固定加点}
\end{chainflow}

请盯住"\textbf{固定}加点"这四个字。陈宁签合同那天谈定的加点，在整个 30 年里\hlmark{一个字都不会改}。会变的只有 LPR 本身，而且要等到他的重定价日才在月供上生效。

拿这套机制去解释一条真实的新闻。

\begin{talkbox}
  \talkline{新闻}{"2026 年 8 月 20 日，1 年期 LPR 报 3.0\%，5 年期以上 LPR 报 3.5\%，均连续多月持平。"}
  \talkline{陈宁}{"不是说货币政策适度宽松吗，怎么一直不降？"}
  \talkline{翻译}{拆成三层看。第一，政策锚没动，7 天期逆回购稳在 1.4\%，报价行的资金成本基准没变。第二，报价行不想动，截至 2026 年二季度末商业银行净息差只有 1.41\%，历史低位，主动下调加点等于自己割自己。第三，负债端是刚的，存款利率降不动，成本降不下来，贷款利率就腾不出空间。\textbf{LPR 按兵不动不是"政策没意愿"，是传导链第 5 环卡住了。}}
\end{talkbox}

一条数据不动，往往比它动了更有信息量，前提是你知道它被链条上哪一环卡住。

\begin{statrow}{3}
  \statitem{1.41}{\%}{商业银行净息差（2026 年二季度末，历史低位）}
  \statitem{3.5}{\%}{5 年期以上 LPR（2026 年 8 月 20 日）}
  \statitem{1.4}{\%}{7 天期逆回购政策利率}
\end{statrow}

\section{那 80 个基点，值 48 万}

"加点"这两个字听着像四舍五入的小数点。把它换算成人民币。

统一条件：300 万元，30 年，等额本息，LPR 为 3.5\%。

\begin{tablewrap}
\begin{xltabular}{\linewidth}{>{\raggedright\arraybackslash}p{0.20\linewidth}Xrrr}
\toprule
\thead{情形} & \thead{执行利率} & \thead{月供} & \thead{30 年总利息} & \thead{与基准之差} \\
\midrule
\textbf{LPR 减 30 个基点} & 3.20\% & 12974 元 & 约 167.1 万 & 每月少 497 元 \\
\textbf{LPR 加 0（陈宁）} & 3.50\% & 13471 元 & 约 185.0 万 & 基准 \\
\textbf{LPR 加 50 个基点} & 4.00\% & 14323 元 & 约 215.6 万 & 每月多 852 元 \\
\bottomrule
\end{xltabular}
\end{tablewrap}

\begin{barfig}
  \barrow{3.20\%（LPR−30bp）}{91}{12974 元/月}
  \barrow{3.50\%（LPR+0）}{94}{13471 元/月}
  \barrow{4.00\%（LPR+50bp）}{100}{14323 元/月}
\end{barfig}

最上面和最下面差 80 个基点。听着是 0.8\%，一个可以忽略的数。

换算成钱：每月 1349 元，三十年 48.5 万元。

而这 80 个基点，往往只取决于签合同那年的市场行情、你的资质，以及你有没有多问过一句"加点还能不能谈"。

\begin{egbox}{举个栗子：LPR 降 25 个基点，陈宁少还多少}
\textbf{情形一，新贷款。}300 万、30 年、等额本息，利率从 3.5\% 降到 3.25\%：月供从 13471 元降到 \textbf{13056 元}，每月少 \textbf{415 元}，全周期少还利息约 \textbf{14.9 万元}。

\textbf{情形二，陈宁这种存量借款人。}他 2024 年 6 月放款，已经还了两年，此时剩余本金约 288.3 万元，剩余期限 28 年。同样降 25 个基点，重定价日之后月供从 13471 元降到约 \textbf{13080 元}，每月少 \textbf{391 元}。

比新贷款少省 24 块。原因是他已经还掉一部分本金，能吃到折扣的基数变小了。\hlmark{降息对存量借款人的帮助，永远小于对新借款人的帮助。}

\end{egbox}

\section{"新闻说降息了，我月供怎么没变"}

这是留言区里出现频率最高的一句抱怨。答案藏在三个词里。

\textbf{重定价日。}浮动利率合同里会约定一个日期，通常是每年 1 月 1 日，或者每年放款日的对应日。只有到了那天，银行才会取用最近一次公布的 LPR 重算月供。8 月降的息，陈宁的重定价日在 1 月 1 日，他要等到明年，中间四个半月照旧按老利率还。

\textbf{固定加点。}LPR 降了，加点不跟着降。当年签的是"LPR 加 80 个基点"，那就永远比基准贵 0.8 个百分点。这是合同，不是政策。

\textbf{批量调整。}也有例外。2023 到 2024 年，中国对存量房贷利率做过集中批量下调，绕开了"加点不能改"这条规则。那是特殊时期的特殊操作，是政策选择，不是常规机制。别把它当成一种可以指望的权利。

\begin{mythbox}{常见误解}
\mvfrow{误解}{"央行一降息，我的房贷马上就便宜了。"}
\mvfrow{事实}{这个说法是错的，中间隔着七个环节。政策利率要先传到银行间市场，再传到 LPR 报价，报价行还得\textbf{愿意}让出自己那 1.41\% 的净息差。就算 LPR 真降了，你还要等自己的重定价日，而加点一分不会动。\hlmark{从政策利率变动到你的月供变动，快则几个月，慢则一年多，也完全可能中途蒸发掉。}}
\end{mythbox}

\section{传导断在哪儿}

上面讲的是"变慢"。更麻烦的是"断掉"。

\begin{flowlist}
  \flowitem{01}{资产负债表衰退}{经济学家辜朝明对日本泡沫破裂后的诊断：资产价格暴跌而债务原封不动，企业和家庭的目标会从"利润最大化"变成"\textbf{负债最小化}"。这时候利率降到 0 也没用，一个净资产为负的人，赚到的每一块钱都会拿去还债。日本走完这个过程用了将近三十年。}
  \flowitem{02}{银行惜贷}{银行资本不足，或者刚吃过一轮坏账，风险偏好整体塌了。它宁可买国债、做同业，也不放给中小企业。\hlmark{钱在银行体系里空转，就是不出门。}}
  \flowitem{03}{没有值得投的项目}{最朴素也最致命的一种。企业算完账发现需求见顶、产能过剩，借钱扩产是纯亏。利率是 3\% 还是 1\%，改不了那张现金流预测表上的答案。}
  \flowitem{04}{利率下限}{名义利率很难显著低于零，因为存款人可以取现金揣着，而现金的名义利率是 0。政策空间到这儿就见底了，剩下的只能靠 QE 和财政。}
\end{flowlist}

中国当下的提前还贷潮，同时踩中了第 1 条和第 3 条：家庭在主动缩表，把手里的存款换成"少欠银行的钱"。按第 7 章的结论，提前还贷会直接注销货币。

于是又是那个别扭的组合：政策在放松，利率在低位，社会感受到的钱依然是紧的。

\begin{pullquote}{}
利率是价格信号，不是命令。价格能改变一个人的意愿，改变不了他"我现在不想借钱"的决心。
\end{pullquote}

顺带说一句机制上的判断，不是建议：提前还贷相当于买入一笔无风险、税后年化收益等于你贷款利率的资产。3.5\% 的房贷，提前还款就是锁定 3.5\% 的无风险回报。这笔账划不划算，取决于你有没有更好的去处、以及那笔钱拿走之后你的现金流还撑不撑得住。这两件事只有你自己知道。

\section{央行控短端，市场定长端}

这条链上最反直觉的一件事留在最后。

央行牢牢控制着曲线最左端，隔夜和 7 天的利率想定多少就是多少。曲线右端的 10 年、30 年，是市场用真金白银投票投出来的，央行只能影响，不能决定。

\begin{egbox}{举个栗子：2026 年 8 月，一件"不该发生"的事}
2026 年 8 月底，美国 10 年期国债收益率升至 \textbf{2025 年 1 月以来的最高水平}。同期美联储的政策利率一动没动，2026 年已连续 5 次维持在 3.50\% 到 3.75\%。

短端纹丝不动，长端往上冲。至少四股力量在同时起作用。

\textbf{①通胀预期。}关税、能源（中东局势）、去全球化都在抬高对未来物价的预期，长期债权人要求更多补偿。\textbf{②期限溢价回归。}QE 时代央行是长债最大的买家，人为压低了这份补偿，央行退场后它自然反弹。\textbf{③财政供给。}美国国债规模 2026 年 8 月突破 \textbf{40 万亿美元}，卖方源源不断，价格自然往下、收益率往上。\textbf{④买家结构在变。}截至 2025 年底，黄金已占全球央行储备 27\%（2024 年底还是 20\%），超越美国国债成为第一大储备资产，美债占比降到 22\%。

沃什本人也承认"名义和实际国债收益率大幅上升"，并指出"从宏观角度看，市场价格表明金融环境已趋紧"。\hlmark{美联储没有收紧，市场替它收紧了。}

\end{egbox}

这件事对普通人的意义很直接。美国的房贷利率主要挂在 10 年期国债收益率上，不是挂在联邦基金利率上。所以"美联储不加息"完全不等于"美国人的房贷不涨"。中国的长期国债收益率同样不由 7 天逆回购单独决定，它反映的是市场对未来十年增长和通胀的整体判断。

\begin{mythbox}{常见误解}
\mvfrow{误解}{"长期利率是央行定的。央行不加息，长期利率就涨不上去。"}
\mvfrow{事实}{央行直接控制的只有曲线最左边那一小段。10 年、30 年的利率由\textbf{通胀预期、期限溢价和债券供求}共同定价。2026 年 8 月的美国就是活标本：政策利率五次按兵不动，10 年期收益率创下 2025 年 1 月以来新高。央行可以影响长端，这正是 QE 存在的理由，但财政供给和通胀预期一起发力的时候，它未必拧得过市场。}
\end{mythbox}

\begin{warnbox}{小心}
本章所有算术都是为了让你看懂机制，不是为了给你一个"该不该提前还贷""该不该换固定利率"的答案。那个答案取决于你的现金流稳定性、有没有更好的资金去处，以及你对未来利率的判断。最后这一项，没有人知道。\textbf{本书不构成投资或借贷建议。}

\end{warnbox}

\bookbreak

\begin{takeaway}{本章一句话}
政策利率到你的月供之间有七道闸，每一道都会削弱、拖慢甚至截断信号。央行控短端，市场定长端。

\begin{itemize}
  \item 链条：政策利率 → 银行间利率（DR007/SOFR）→ 同业存单与短债 → 收益率曲线 → 银行负债成本 → LPR 加点 → 你的月供。
  \item 2026 年 8 月 LPR 连续多月不动，卡点在 1.41\% 的净息差和降不下来的存款成本，不是政策没意愿。
  \item 300 万 30 年，加点从减 30bp 到加 50bp，月供差 1349 元，总利息差 48.5 万元。
  \item 传导会断在四个地方：资产负债表衰退、银行惜贷、没有值得投的项目、利率下限。
\end{itemize}

\end{takeaway}

\begin{bridgebox}
\textbf{下一章}说了半天利率，它到底长在哪个市场上？答案是一个规模比全球股市大得多、每天决定着陈宁的月供和一国财政命运、却几乎没人在饭桌上聊起的地方。
\end{bridgebox}

